\documentclass[11pt]{article}
\usepackage[T1]{fontenc}
\usepackage{lmodern}
\usepackage[letterpaper,margin=1in]{geometry}
\usepackage{amsmath,amssymb,amsthm,booktabs}
\usepackage[colorlinks=true,linkcolor=black,citecolor=black,urlcolor=blue]{hyperref}
\hypersetup{pdftitle={Cyclic substitutions and a spectral short-cycle conjecture},
pdfauthor={Alex Chan, Weilin Meng, Kurt Sweat}}
\newtheorem{theorem}{Theorem}
\newtheorem{lemma}[theorem]{Lemma}
\newtheorem{corollary}[theorem]{Corollary}
\newcommand{\rhoA}{\rho(A)}
\makeatletter
\renewcommand{\@maketitle}{\begin{center}{\Large\bfseries\@title\par}\vskip0.8em{\normalsize\@author\par}\vskip0.6em{\normalsize\@date\par}\end{center}}
\makeatother
\title{Cyclic substitutions and a spectral short-cycle conjecture}
\author{Alex Chan$^{1}$\quad Weilin Meng$^{2}$\quad Kurt Sweat$^{3}$\\[6pt]
  \small $^{1}$Harvard University and National Bureau of Economic Research\\
  \small $^{2}$Harvard Business School\\
  \small $^{3}$UT Southwestern Medical Center\\[3pt]
  \small \texttt{achan@hbs.edu}}
\date{31 October 2026}

\begin{document}
\maketitle

\begin{abstract}
The spectral short-cycle statement attributed to Charbit in Sullivan's survey asserts that an $n$-vertex loopless digraph with adjacency spectral radius at least $n/k$ has a directed cycle of length at most $k$. We show that this statement holds for positive integers $k$ exactly when $k\leq2$. For the negative cases, we apply a previously studied construction obtained by substituting transitive tournaments into a directed cycle. An explicit positive eigenvector gives its spectral radius and yields counterexamples for every $k\geq3$. The cases $k=3,4,5$ have witnesses on twenty, fifteen and twelve vertices, respectively. The proofs are self-contained apart from standard Perron--Frobenius facts. A companion Lean formalization verifies the twelve-vertex counterexample, including its strict spectral inequality.
\end{abstract}

\section{The statement and the counterexamples}

For a finite simple loopless digraph $D$ on $n\geq1$ vertices, let $A=A(D)$ be its zero--one adjacency matrix, with $A_{uv}=1$ precisely when $u\to v$. Write $g(D)$ for its directed girth, with $g(D)=\infty$ if $D$ is acyclic, and $\rho(A)$ for the maximum modulus of its complex eigenvalues. For positive integer $k$, the statement recorded as Conjecture~6.18 of Sullivan~\cite{sullivan}, attributed there to Charbit, is
\begin{equation}\label{eq:statement}
 \rho(A)\geq\frac{n}{k}\quad\Longrightarrow\quad g(D)\leq k.
\end{equation}
We address this formulation exactly as printed. In particular, it imposes no regularity or Eulerian hypothesis. An earlier spectral formulation appears in Charbit's thesis~\cite[Section~5.4]{charbit}.

The construction used below is already present in the literature. Chudnovsky, Seymour and Sullivan~\cite[Section~1]{css} give $C_4[T_s]$ as an extremal example for a feedback-arc-set problem. Aharoni, Berger, Chudnovsky, Guo and Zerbib~\cite[Section~2]{aharoni} describe the general family $C_g[T_s]$, crediting Hompe's 2019 thesis, and use it to disprove a proposed upper bound on directed girth in terms of reciprocal outdegrees. We apply this established family to~\eqref{eq:statement} by evaluating its adjacency spectral radius. The graph family and the substitution operation are not claimed as new.

\begin{theorem}\label{thm:main}
For a positive integer $k$, implication~\eqref{eq:statement} holds for all finite simple loopless digraphs on at least one vertex if and only if $k\leq2$. For every $k\geq3$ there is a simple, strongly connected oriented graph $D$ such that $g(D)=k+1$ and $\rho(A(D))>|V(D)|/k$.
\end{theorem}

\section{The two positive cases}

For $k=1$, every row sum is at most $n-1$, so $\rho(A)\leq n-1<n$ and the antecedent is impossible.

For $k=2$, suppose that $D$ has no digon. Then $A+A^{\mathsf T}\leq J-I$ entrywise, where $J$ is the all-ones matrix. A nonnegative matrix, including a reducible one, has a nonzero nonnegative real vector $x$ with $Ax=\rho(A)x$; see, for example, Butler--Siegel~\cite[Remark~3.1]{butler}. As $x^{\mathsf T}x>0$, Cauchy--Schwarz gives
\[
 2\rho(A)x^{\mathsf T}x
 =x^{\mathsf T}(A+A^{\mathsf T})x
 \leq\left(\sum_i x_i\right)^2-x^{\mathsf T}x
 \leq(n-1)x^{\mathsf T}x.
\]
Thus $\rho(A)\leq(n-1)/2<n/2$. These standard arguments establish the positive part of Theorem~\ref{thm:main}, including $n=1$.

\section{Cyclic substitution}

Let $g\geq3$ and $s\geq1$ be integers. Define $D_{g,s}$ on
\[
 V(D_{g,s})=\mathbb Z/g\mathbb Z\times\{0,\ldots,s-1\}
\]
by putting in exactly the arcs
\begin{equation}\label{eq:arcs}
 (i,a)\to(i,b)\quad(a<b),\qquad
 (i,a)\to(i+1,b)\quad(0\leq a,b<s).
\end{equation}
Equivalently, $D_{g,s}=C_g[T_s]$, with $T_s$ the transitive tournament of order $s$. We use the construction described in~\cite{css,aharoni}.

\begin{lemma}\label{lem:family}
The digraph $D_{g,s}$ is simple, oriented and strongly connected. Moreover,
\[
 |V(D_{g,s})|=gs,\qquad
 |E(D_{g,s})|=g\left(s^2+\binom{s}{2}\right),\qquad
 g(D_{g,s})=g,
\]
and
\begin{equation}\label{eq:radius}
 \rho(A(D_{g,s}))=\frac{1}{2^{1/s}-1}.
\end{equation}
\end{lemma}

\begin{proof}
The counts follow from~\eqref{eq:arcs}. There are no loops or opposite arcs, since $g\geq3$ and the arcs within a class follow a strict order. Crossing arcs allow every vertex to reach every other vertex, going once around the classes if necessary.

A directed cycle cannot use only within-class arcs. The positive number of crossing arcs in a cycle is divisible by $g$, since each crossing adds one to the class index modulo $g$. Its length is therefore at least $g$. Keeping the second coordinate zero gives a cycle of length $g$.

Put $r=2^{1/s}>1$, $\lambda=(r-1)^{-1}$, and $x(i,a)=r^{s-1-a}$. Since $r^s=2$,
\[
 S:=\sum_{b=0}^{s-1}r^{s-1-b}
 =\frac{r^s-1}{r-1}=\lambda.
\]
The sum of the coordinates on out-neighbors of $(i,a)$ is
\[
 (Ax)(i,a)=S+\sum_{b=a+1}^{s-1}r^{s-1-b}
 =\lambda+\frac{r^{s-1-a}-1}{r-1}
 =\lambda x(i,a).
\]
Thus $\lambda$ is an eigenvalue with a strictly positive eigenvector.
For completeness, let $Az=\mu z$ with $z\neq0$ complex, and choose $u$ maximizing $|z_u|/x_u$. This maximum is positive, and
\[
 |\mu|\,|z_u|\leq\sum_v A_{uv}|z_v|
 \leq\frac{|z_u|}{x_u}\sum_v A_{uv}x_v
 =\lambda|z_u|.
\]
Hence every eigenvalue satisfies $|\mu|\leq\lambda$, which proves~\eqref{eq:radius}.
\end{proof}

\begin{proof}[Proof of the negative part of Theorem~\ref{thm:main}]
Use $D_{k+1,5}$. It has $n=5(k+1)$ vertices and directed girth $k+1$. The integer inequality
\[
 23^5=6\,436\,343>6\,400\,000=2\cdot20^5
\]
gives $2^{1/5}<23/20$ and hence
\[
 \rho(A)=\frac{1}{2^{1/5}-1}>\frac{20}{3}
 \geq\frac{5(k+1)}{k}=\frac nk,
\]
where the second inequality uses $k\geq3$.
\end{proof}

\section{Small witnesses and exact certificates}

Smaller tournament sizes suffice at $k=4$ and at every $k\geq5$. For $k=4$, use $D_{5,3}$: the inequality $19^3>2\cdot15^3$ gives $\rho(A)>15/4$. For $k\geq5$, use $D_{k+1,2}$, since
\[
 \rho(A)=1+\sqrt2>\frac{12}{5}
 \geq\frac{2(k+1)}{k}.
\]
Thus three small instances are as follows; their orders are not claimed to be minimal.
\begin{center}
\begin{tabular}{ccrrcc}
\toprule
$k$ & Digraph & Vertices & Arcs & Girth & Spectral radius\\
\midrule
3 & $D_{4,5}$ &20&140&4&$(2^{1/5}-1)^{-1}>20/3$\\
4 & $D_{5,3}$ &15&60&5&$(2^{1/3}-1)^{-1}>15/4$\\
5 & $D_{6,2}$ &12&30&6&$1+\sqrt2>12/5$\\
\bottomrule
\end{tabular}
\end{center}

Short integer certificates give another proof of every negative case. Assign the same positive vector $w$ to each class. Direct sums of outgoing weights give the following vectors and residuals, displayed per class:
\[
\begin{array}{ccl}
k=3:&w=(33,29,25,22,19),&3Aw-20w=(9,2,7,1,4),\\[2pt]
k=4:&w=(11,9,7),&4Aw-15w=(7,1,3),\\[2pt]
k\geq5:&w=(17,12),&kAw-2(k+1)w=(7k-34,5k-24).
\end{array}
\]
In the last row, the residuals equal $7(k-5)+1$ and $5(k-5)+1$, so every residual is strictly positive. The standard positive-vector bound gives $\rho(A)>n/k$; see~\cite[Theorem~1.2]{butler}. For completeness, if $y\geq0$ is a nonzero left Perron vector, then $y^{\mathsf T}w>0$, and $Aw>cw$ gives
\[
 \rho(A)y^{\mathsf T}w=y^{\mathsf T}Aw>c\,y^{\mathsf T}w.
\]
This also applies to reducible nonnegative matrices. These integer certificates supplement the closed-form proof above.

The minimum outdegree of $D_{g,s}$ is $s$, attained at $(i,s-1)$. In our examples $g=k+1$, so $s<gs/k$. Consequently they do not satisfy the degree hypothesis in the ordinary Caccetta--H\"aggkvist conjecture.

\section{Formal verification of the twelve-vertex example}

The twelve-vertex example also has a Lean~4 formalization, maintained with the companion sources at
\begin{center}
\url{https://github.com/alexgocardinal/charbit-spectral-counterexamples}.
\end{center}
The development uses Lean~4.24.0 and Mathlib at tag \texttt{v4.24.0}. In the source version \texttt{ee815b3}, the theorem \texttt{Charbit.counterexample} verifies looplessness, the absence of opposite arcs, the absence of nonempty closed directed walks of lengths one through five, and the strict inequality $\rho(A)>12/5$. Here the spectral radius is Mathlib's standard supremum of the moduli of the complex spectrum, applied to the complexified zero--one adjacency matrix. The proof verifies the eigenvalue $1+\sqrt2$ using the repeated vector $(\sqrt2,1)$ and then bounds the spectral radius from below by its modulus. The arcs of the displayed six-cycle are checked separately.\footnote{The source files are \texttt{Charbit.lean} and \texttt{Verify.lean}. The commands \texttt{lake build} and \texttt{lake env lean Verify.lean} reproduce the build and assumption audit. The counterexample theorem depends only on Lean's standard foundational axioms \texttt{propext}, \texttt{Classical.choice} and \texttt{Quot.sound}; there are no admitted proof steps. The repository is private during manuscript preparation.}

This formalization covers the twelve-vertex counterexample. The variable-order family and the classification in Theorem~\ref{thm:main} are proved in the preceding sections and are not asserted to have been formalized.

\begin{thebibliography}{9}

\bibitem{aharoni}
R. Aharoni, E. Berger, M. Chudnovsky, H. Guo and S. Zerbib,
Nonuniform degrees and rainbow versions of the Caccetta--H\"aggkvist conjecture,
\emph{SIAM J. Discrete Math.} \textbf{37} (2023), 1704--1714.
\url{https://doi.org/10.1137/22M1529658}.

\bibitem{butler}
B. K. Butler and P. H. Siegel,
Generalized sharp bounds on the spectral radius of digraphs,
\emph{Linear Algebra Appl.} \textbf{439} (2013), 1468--1478.
\url{https://doi.org/10.1016/j.laa.2013.04.029}.

\bibitem{charbit}
P. Charbit, \emph{Plongements de graphes et \'etude des circuits},
doctoral thesis, Universit\'e Claude Bernard--Lyon I, 2005, Section~5.4.
\url{https://www.irif.fr/~charbit/Recherche/publis/these_fr.pdf}.

\bibitem{css}
M. Chudnovsky, P. Seymour and B. Sullivan,
Cycles in dense digraphs, \emph{Combinatorica} \textbf{28} (2008), 1--18.
\url{https://doi.org/10.1007/s00493-008-2331-z}.

\bibitem{sullivan}
B. D. Sullivan,
\emph{A Summary of Results and Problems Related to the Caccetta--H\"aggkvist Conjecture},
AIM 2006--13, arXiv:math/0605646v1, 2006, Conjecture~6.18.
\url{https://arxiv.org/abs/math/0605646}.

\end{thebibliography}
\end{document}
